\FloatBarrier
\section{Discussion}
Our formulation makes clear which part of the result comes from the representation,
the regularization, the optimizer and the available information. Exporting a fit
on 256 bins reduces display quantization; it does not create 256 independent
diffusion observations. Likewise, a lower prediction error cannot by itself
identify distributional width, chemical species or a unique factorization.

The support selector limits the assignment of a proposed rate to unrelated
spectral regions. It first asks whether that rate contributes beyond the other
candidate decays, using groups of channels to retain coherent weak signals.
Unrestricted support remains a candidate. The factorial study finds contributions
from both screening and predictive selection, together with an interaction.
It does not establish an independent advantage for the RAI or DOME proposals.
Their primary outputs coincide, as do many of the additional outputs, although
the weak-component panel contains one different minority-rate match. The searches
have not been proved equivalent.

The spectral groups are geometric windows and need not coincide with chemical
resonances. Masking can remove weak signal, errors in nuisance rates can bias the
residualized score, and nearly coincident rates make its direction ill-conditioned.
Proposition~\ref{prop:screen} and the stability bounds describe fixed designs.
The 300 null matrices and 96 weak-component problems characterize particular
generators; they do not calibrate $\alpha=0.01$ as a universal family-wise error
rate. The prediction allowance uses the correct difference between fitted
predictions, but choosing the best model and searching supports make its subsequent
use heuristic. The discovery masks in the new weak-component and experimental
studies exclude validation observations. This removes one source of data reuse,
while leaving uncertainty in the estimated nuisance model and the need for
selection-adjusted inference.

A TRAIn-MF factor represents a shared distribution, which may be broad or
multimodal. Its unit-mass normalization fixes scale, not chemical identity.
The runtime curvature check concerns injectivity of a conditional subproblem.
The controlled warm-start study supplies factor count to isolate optimization
and records nonstationary outcomes. It therefore tests neither automatic species
selection nor convergence to a global minimum. The search caps and iteration
budgets remain essential parts of the reported procedure.

An atomic model can predict the measured attenuation of a broad distribution
well because the Laplace operator smooths differences between measures.
The bound in Section~7.3 explains why good prediction can coexist with an
incorrect distributional model. Screening may reduce estimation variance at
the cost of approximation bias, but the prediction gains alone do not isolate
that mechanism or recover the true width. This matters for industrial resins,
whose heterogeneity may require more than a few atoms. Their selected model
order should not be read as a count of molecular constituents.

The external comparisons show how strongly the result depends on the fitted
representation and the chosen evaluation measure. Signed cancellation, native
order rules and incomplete optimization remain visible alongside prediction
error. A signed fit may not define a positive measure, so transport scores
computed only on eligible outputs describe different subsets. On the public
mixture, signed SCORE converges and passes the stability checks; it and LOCODOSY
predict better than the positive global estimators. We report that advantage
without treating it as proof of chemical recovery. The nominal composition
provides neither a complete reference order nor independently measured diffusion
coefficients. These results also cannot rank the four external methods that
were not run.

The larger neural study compares five seeded policies on 120 new problems,
including changes in noise and acquisition design. All policies use the same
search caps and numerical acceptance checks. The policies share a trace bank
and test problems; their 720 fitted outputs are not 720 independent mixtures.
We average seed-level metrics within each problem before the paired analysis
and retain numerical failures. The pooled accuracy intervals include zero. The cost
reduction is conditional on the stated training and search design and does
not establish an end-to-end speedup or experimental chemical recovery.

The resin measurements provide four specimen series, each with three
diffusion-time blocks. They do not provide twelve independently prepared
mixtures of known composition. Their strongly correlated background and possible
phase, baseline, suppression, exchange, relaxation or encoding errors can violate
the ideal observation model. Support is detected without the reserved acquisitions;
negative coherent structure and unestimated regions remain documented.
Prediction on reserved acquisitions and consistency across diffusion times are
useful checks, but neither supplies a component diffusion reference or a
molecular-weight calibration. The experimental interpretation must remain within
those limits.

\section{Conclusions}
Joint positive Laplace inversion can improve spectral assignment when shared
components are supported by the data. The distinction between bin masses,
densities, atoms and shared profiles is essential to interpreting the result.
Our conditional stability and uniqueness calculations explain where reconstruction
can fail; the additional experiments show those limits in practice. Weak rates
are missed, finite optimization budgets leave some fits nonstationary, and the
neural accuracy gains remain inconclusive. The observed reduction in spectral
leakage therefore supports a specific empirical contribution. Chemical identity
and distributional width require additional evidence.